\bookchapter{Capitolo 28: Spendere bene per il telefono desiderato — Guida definitiva all'acquisto}

\chapterintro{Per l'acquisto di un telefono, molti guardano recensioni per una settimana, confrontano i prezzi per un mese, e poi finiscono per spendere cinque-seimila yuan in un top di gamma che in realtà non sfruttano. Questo capitolo non elenca schede tecniche — quelle le trovi sul sito ufficiale. Voglio darti un “metodo di scelta”, così in futuro, qualsiasi telefono nuovo esca, potrai giudicare da solo se è adatto a te.}

\booksection{Metodo in tre passi: esigenze budget modello}

Molti online chiedono “quale telefono comprare con tremila yuan”, domanda che in realtà non ha risposta — perché non si sa cosa ci vuoi fare. L'ordine corretto di pensiero quando si acquista un telefono è:

\begin{tcbinfo}{I tre colpi della scelta}
\begin{enumerate}
  \item \textbf{Prima pensa a cosa ti importa di più}: foto? giochi? autonomia? schermo? fluidità del sistema? o ti basta poter chiamare e usare WeChat? Metti in ordine le tue esigenze.
  \item \textbf{Poi fissa un tetto di budget}: non "circa quanto", ma "al massimo quanto posso spendere". Fissa un limite rigido, così non guardi modelli sempre più cari.
  \item \textbf{Solo alla fine scegli il modello}: con le tue priorità e il budget, cerca il telefono corrispondente.
\end{enumerate}

\end{tcbinfo}

\booksection{Come leggere le specifiche principali — non serve capirle tutte, cinque bastano}

Le schede tecniche hanno decine di voci, ma l'utente comune deve guardare solo cinque dimensioni:

\begin{center}
\begin{tabularx}{\linewidth}{|X|X|X|}
\hline
\textbf{Dimensione} & \textbf{Come valutarla} & \textbf{In una frase} \tabularnewline
\hline
Processore (SoC) & Determina quanti anni il telefono resterà fluido. Standard 2026: Snapdragon 8 Elite Gen5/Dimensity 9500 o superiore = top; Snapdragon 8s Gen4/Dimensity 8400 = fascia medio-alta; Snapdragon 7+ Gen4/Dimensity 7350 = sufficiente per tre anni. Sotto la serie Snapdragon 6, valuta con cautela. & Compra il nuovo non il vecchio; il SoC determina quanti anni il telefono resterà fluido \tabularnewline
Schermo & OLED ha colori migliori e consuma meno rispetto a LCD; refresh rate 120Hz scorre molto meglio di 60Hz; risoluzione almeno 1080P o superiore. PWM ad alta frequenza (oltre 1920Hz) è più amichevole per gli occhi. & OLED + 120Hz + PWM ad alta frequenza = soglia minima per un buon schermo \tabularnewline
Fotocamera & \textbf{Non guardare i megapixel, guarda la dimensione del sensore}. Un sensore grande batte uno piccolo; la dimensione del sensore conta molto più dei megapixel. Il grandangolo è utile; il tele (zoom ottico 3x o più) per oggetti lontani; il macro spesso è riempitivo. & Sensore principale grande + tele utile > tante lenti riempitive \tabularnewline
Batteria e ricarica & 4500mAh è la soglia minima, sopra 5000mAh si fa una carica al giorno con tranquillità. Per la ricarica rapida guarda i watt: oltre 60W è veloce, oltre 100W è molto veloce. La ricarica wireless è comoda ma lenta. & 5000mAh + 60W = combo senza ansia \tabularnewline
Impegno di aggiornamento sistema & Determina per quanto il telefono resterà sicuro. Apple in genere garantisce 5-6 anni di aggiornamenti; i top di gamma Samsung 4 anni sistema + 5 anni aggiornamenti di sicurezza; Google Pixel 7 anni; i cinesi di fascia media hanno impegni più brevi (2-3 anni). & Se vuoi tenerlo a lungo, scegli brand con impegni lunghi di aggiornamento \tabularnewline
\hline
\end{tabularx}
\end{center}

\booksection{Raccomandazioni per fascia di prezzo (riferimento al mercato cinese a metà 2026)}

Le raccomandazioni seguenti si basano sul mercato cinese continentale a metà 2026, i prezzi sono quelli di partenza ufficiali. I prezzi effettivi possono variare, fa fede quello al momento dell'acquisto.

\booksubsection{Sotto i 1000 yuan: può chiamare, può usare WeChat, basta che funzioni}

\textbf{A chi si addice}: studenti, telefono di riserva, esigenze modeste, uso principale per chiamate e WeChat.

\textbf{Cosa puoi trovare}: Redmi Note serie, Honor X serie, realme numero serie nelle versioni medio-basse. Processore circa Snapdragon serie 6 o Dimensity 6000, schermo prevalentemente LCD, foto nitide ma autofocus lento, fluidità di sistema nella media.

\textbf{Consigli d'acquisto}: in questa fascia guarda due cose: la batteria (scegli 5000mAh o più) e la memoria (almeno 6GB+128GB, evita i 4GB che si bloccano). Redmi in questa fascia è il più competitivo.

\booksubsection{1000-2000 yuan: il miglior rapporto qualità-prezzo, la zona “dolce” dove si usa 3-4 anni senza lag}

\textbf{A chi si addice}: studenti, budget limitato ma voglia di durata, senza esigenze estreme su foto e giochi.

\textbf{Modelli rappresentativi}: Redmi Turbo serie, iQOO Z serie, Honor X serie alta, OnePlus Ace serie entry.

\textbf{Consigli d'acquisto}: scegli modelli con processore Snapdragon 7+ Gen3 o Dimensity 7300 o superiore, con 8GB+256GB. Questa configurazione per uso quotidiano + partita occasionale a Honor of Kings è più che sufficiente, dura 3-4 anni senza lag.

\booksubsection{2000-3500 yuan: la fascia centrale, prestazioni e foto migliorano nettamente}

\textbf{A chi si addice}: utenti comuni che vogliono buone foto, giochi fluidi e durata. È la fascia più comprata.

\textbf{Modelli rappresentativi}: Xiaomi serie numero (es. Xiaomi 17 standard dopo lo sconto), vivo S serie alta, OPPO Reno serie, Honor serie numero, Redmi K serie.

\textbf{Consigli d'acquisto}: la caratteristica di questa fascia è il “telefono tuttofare” — senza punti deboli evidenti. Rispetto ai modelli sotto i 2000 yuan, i miglioramenti principali sono: fotocamera nettamente migliore (con stabilizzazione ottica), schermo OLED di alta qualità, ricarica più veloce (spesso oltre 80W), qualità costruttiva più raffinata. Scegli 12GB+256GB o 12GB+512GB, dura 4-5 anni senza problemi.

\booksubsection{3500-5000 yuan: quasi top di gamma, per tutti tranne gli utenti estremi}

\textbf{A chi si addice}: utenti esigenti sulla fotocamera, chi vuole qualità costruttiva top.

\textbf{Modelli rappresentativi}: Xiaomi serie numero standard, vivo X serie, OPPO Find X serie standard, iPhone top di gamma dell'anno precedente (es. iPhone 16).

\textbf{Consigli d'acquisto}: in questa fascia puoi già trovare i top di gamma dell'anno precedente. Il mio consiglio è:\textbf{compra il top di gamma dell'anno scorso, non il medio di gamma di quest'anno}. Materiali, costruzione, hardware fotografico del top di gamma dell'anno precedente sono in genere un gradino sopra il medio di gamma attuale. A patto che tu accetti che non sia “l'ultimissimo modello”.

\booksubsection{Oltre 5000 yuan: fascia top, per chi non ha limiti di budget o cerca il massimo in una specifica esperienza}

\textbf{A chi si addice}: appassionati di fotografia, giocatori incalliti, chi cerca schermo e lavorazione top, creatori di contenuti video con il telefono.

\textbf{Modelli rappresentativi}: iPhone 17 Pro/Pro Max, Samsung S26 Ultra, Huawei Pura 80 Ultra, Xiaomi 17 Ultra, vivo X300 Pro.

\textbf{Consigli d'acquisto}: in questa fascia le differenze tra i brand sono sempre più piccole, in realtà conta soprattutto quale sistema preferisci. Se cerchi foto scegli vivo o Huawei, per i video iPhone, per schermo e S-Pen Samsung, per un equilibrio Xiaomi. A dire il vero a questa cifra,\textbf{puoi seguire il cuore, scegli quello che ti piace guardare}.

\booksection{Scegliere in base allo scenario — qual è la tua esigenza?}

Ecco alcuni dei “profili d'uso” più comuni, basta trovare il proprio:

\begin{center}
\begin{tabularx}{\linewidth}{|X|X|X|}
\hline
\textbf{Scenario d'uso} & \textbf{Cosa guardare} & \textbf{Direzione consigliata} \tabularnewline
\hline
Appassionato di foto & Sensore principale grande (1 pollice ideale), stabilizzazione ottica, tele (zoom ottico 3x o più), chip imaging & vivo X serie, Huawei Pura/Mate serie, Xiaomi Ultra serie \tabularnewline
Gamer & Processore top, sample rate touch alto, sistema di dissipazione, batteria grande + ricarica rapida & iQOO serie numero, ROG Phone, Red Magic, Redmi K serie \tabularnewline
Autonomia prima di tutto & Batteria oltre 5500mAh, processore efficiente, schermo LTPO & Honor Magic serie, OnePlus, Redmi Note versione grande batteria \tabularnewline
Per gli anziani & Modalità semplice, caratteri grandi, volume alto, batteria grande, molti centri assistenza & Huawei Enjoy, Honor Play, OPPO A serie, Redmi Note \tabularnewline
Per studenti & Rapporto qualità-prezzo, durevole, batteria grande, controllo genitori & Redmi Note/Turbo, realme GT, iQOO Neo \tabularnewline
Utente ecosistema Apple & Ha già Mac/iPad/AirPods, iPhone è l'unica soluzione & iPhone 17 standard (rapporto qualità-prezzo) o Pro (esperienza completa) \tabularnewline
Viaggi all'estero & Supporto eSIM, copertura completa delle bande globali, servizi Google & iPhone, Samsung S serie, Pixel, OnePlus \tabularnewline
\hline
\end{tabularx}
\end{center}

\booksection{Cinque trappole da evitare — controlla prima di comprare}

\begin{enumerate}
  \item \textbf{Trappola 1: quattro o cinque fotocamere non significano buone foto}— ai produttori piace sommare lenti nelle specifiche, spesso due sono da 2 megapixel "riempitive" (macro + profondità), in realtà inutili. Per valutare la fotocamera,\textbf{guarda solo sensore principale + grandangolo + tele}. Le altre puoi ignorarle.
  \item \textbf{Trappola 2: non farti incantare da 100 megapixel}— tanti megapixel non significano buona qualità. Dimensione del sensore, apertura, algoritmi sono molto più importanti dei megapixel. Molti top di gamma con sensore grande da 50MP fanno foto molto migliori di alcuni telefoni da 1000 yuan con 100MP.
  \item \textbf{Trappola 3: non comprare telefoni brandizzati dall'operatore}— "ricarica e ti diamo il telefono" sembra un affare, ma i modelli brandizzati hanno aggiornamenti lenti, App preinstallate non rimovibili, e spesso sono sim-locked (utilizzabili solo con quella SIM). Sembra un risparmio, ma in realtà ti "legano".
  \item \textbf{Trappola 4: "schermo curvo = fascia alta" è marketing, non verità}— gli schermi curvi erano un tempo segno di top di gamma, ma sempre più utenti ne scoprono i problemi: pellicole difficili, tocchi accidentali sui bordi, riparazione costosa se si rompe. La tendenza 2026 è tornare al "micro-curvo" o "totalmente piatto". Non scegliere uno schermo curvo solo perché "è più nobile".
  \item \textbf{Trappola 5: non tutti i "telefoni IA" sono davvero utili}— nel 2026 tutti i produttori gridano al "telefono IA", ma molti si limitano a rinominare l'assistente vocale esistente. Le funzioni IA davvero utili sono: traduzione in tempo reale, rimozione oggetti IA (eliminare persone dalle foto), fotoritocco IA, inferenza on-device di grandi modelli, riassunti intelligenti localizzati. Le altre funzioni vistose sono per lo più slogan.
\end{enumerate}

\booksection{Quando si compra al miglior prezzo?}

\begin{itemize}
  \item \textbf{618 (fine maggio - metà giugno)}: il periodo di maggiori sconti dell'anno. Per spingere le vendite, i top di gamma scendono di 800-1500 yuan. È la finestra migliore per acquisti nella prima metà dell'anno.
  \item \textbf{Double 11 (fine ottobre - metà novembre)}: promozioni di seconda metà dell'anno, simili al 618. In questo periodo molte marche lanciano la nuova generazione di top di gamma, e i modelli precedenti scendono di più.
  \item \textbf{Un mese dopo il lancio di un nuovo modello}: al lancio prezzo pieno, dopo un mese i canali terzi scendono generalmente di 200-500 yuan.
  \item \textbf{Da due settimane prima del Capodanno cinese alla Festa delle Lanterne}: i negozi fisici svuotano le scorte con la "festa dei doni di Capodanno", puoi contrattare con il commesso.
  \item \textbf{Sconsigliato}: comprare al lancio il giorno stesso — a meno che tu non voglia proprio essere tra i primi utenti, il prezzo di lancio è il più alto.
\end{itemize}

\booksection{Quale canale scegliere — dove conviene comprare?}

\begin{itemize}
  \item \textbf{Brand store ufficiali / flagship Tmall / JD self-operated}: i più sicuri, prezzi non aggressivi ma nessun rischio di prodotti ricondizionati o falsi, post-vendita migliore.
  \item \textbf{"Sussidio da 10 miliardi" di Pinduoduo}: prezzi più bassi, ma scegli negozi con etichetta "sussidio da 10 miliardi". Senza etichetta, non toccarli. Ricorda di filmare l'unboxing, in caso di problemi puoi restituirli.
  \item \textbf{Negozi fisici}: puoi toccare e provare la sensazione in mano. Attenzione: i commessi propongono i modelli con margine più alto (non i più adatti a te), conviene scegliere 2-3 modelli target online prima di entrare, e dire "voglio vedere solo questi".
  \item \textbf{Piattaforme di seconda mano (Xianyu / Zhuanzhuan)}: più economici ma rischiosi. Se compri usato, ricorda tre cose: chiedi fattura o screenshot dell'ordine, chiedi l'IMEI per verificarlo,\textbf{usa solo la transazione tramite piattaforma, non bonifico privato}.
\end{itemize}

\begin{tcbpractice}{Fai: usa questo metodo per scegliere un telefono}
\begin{enumerate}
  \item Prendi un foglio e scrivi le due esigenze principali (es. autonomia prima, foto seconda).
  \item Fissa un tetto di budget (es. 3500 yuan).
  \item Apri JD o Tmall, filtra per prezzo, individua 3-5 modelli candidati.
  \item Cerca su Bilibili "modello XX, esperienza dopo un mese" (non recensioni del primo giorno, ma impressioni reali dopo almeno un mese), guarda due o tre video.
  \item Conferma l'affidabilità del canale e ordina; all'arrivo, riprendi l'unboxing.
\end{enumerate}

\end{tcbpractice}

\begin{tcbsummary}{Riepilogo del capitolo}
\begin{itemize}
  \item Tre colpi della scelta: prima le esigenze poi il budget infine il modello
  \item Specifiche chiave solo cinque: processore, schermo, fotocamera, batteria/ricarica, aggiornamenti di sistema
  \item 1500-3500 yuan è la fascia "dolce" con il miglior rapporto qualità-prezzo
  \item Comprare il top di gamma dell'anno scorso è più conveniente del medio di gamma di quest'anno
  \item Cinque trappole: truffa delle multi-fotocamere, trappola degli alti megapixel, brandizzazione operatore, mito dello schermo curvo, slogan IA
  \item 618 e Double 11 sono le due finestre migliori dell'anno per comprare
\end{itemize}

\end{tcbsummary}
